\subsection{PROPERTIES OF FINITE CHARACTER}

DEFINITION 2. Let E be a set. A set S of subsets of E is said to be of finite character if the relation X ∈ S is equivalent to the relation ``every finite subset of X belongs to S''.

A property \(P\{X\}\) of a subset X of a set E is said to be of finite character if the set of subsets X of E for which \(P\{X\}\) is true is of finite character.

\minorheading{Examples}

\begin{enumerate}
\def\labelenumi{(\arabic{enumi})}
\tightlist
\item
  The set of totally ordered subsets of an ordered set E is of finite character. Indeed, a subset X of E is totally ordered if and only if every subset of X consisting of two elements is totally ordered.
\end{enumerate}

* (2) The set of all free subsets of a module is of finite character. The same is true of the set of all algebraically free subsets of an extension of a field.

\begin{enumerate}
\def\labelenumi{(\arabic{enumi})}
\setcounter{enumi}{2}
\tightlist
\item
  The set of submodules of a module E is not of finite character, because a finite subset of a submodule of E is not necessarily a submodule of E. *
\end{enumerate}

THEOREM 1. Every set S of subsets of a set E which is of finite character has a maximal element (when ordered by inclusion).

By Theorem 2 of §2, no. 4 it is enough to show that \(\mathfrak{S}\) is inductive. To do this we shall show that if \(\mathfrak{G}\) is any subset of \(\mathfrak{S}\) which is totally ordered by inclusion, the union X of the sets of \(\mathfrak{G}\) belongs to \(\mathfrak{S}\) (§2, no. 4, Theorem 2, Corollary 2). Since \(\mathfrak{S}\) is of finite character, it is enough to show that every finite subset Y of X belongs to \(\mathfrak{S}\) . Now, for each \(y \in Y\) , there exists a set \(Z_y \in \mathfrak{G}\) such that \(y \in Z_y\) . Since the set of sets \(Z_y\) ( \(y \in Y\) ) is finite and totally ordered by inclusion, it has a greatest element S (no. 4, Corollary 1 to Proposition 3); in other words, there exists a set \(S \in \mathfrak{G}\) such that \(Y \subset S\) . But since \(S \in \mathfrak{S}\) and since Y is a finite subset of S, we have \(Y \in \mathfrak{S}\) because \(\mathfrak{S}\) is of finite character; this completes the proof.

